\documentclass[italian,12pt]{article}

\def\Title{Verbale interno}
\def\Subtitle{Riunione interna settimanale}
\def\Author{7Last}
\def\Date{2024-05-22}
\def\Version{v1.0}

%pacchetti extra da scaricare dblfloatfix, fancyhdr
\usepackage[left=2cm, right=2cm, bottom=3cm, top=3cm]{geometry}
\usepackage{fancyhdr}%creazione header-footer
\usepackage{graphicx} %serve per inserire immagini
\graphicspath{ {../../../logo/} }
%\usepackage{dblfloatfix} %serve per posizionare gli elementi dove si vuole
\usepackage[hidelinks]{hyperref} %serve per i link
\usepackage{tikz}
\usepackage{tgadventor} % font
\usepackage[useregional=numeric,showseconds=true,showzone=false]{datetime2}
\usepackage{caption}

\usepackage{hyperref}
\usepackage{tocloft}
\usepackage{titlesec}
\usepackage{color}
\usepackage{ulem}
\usepackage{pgfplots}
\usepackage{pgf-pie}
\usepackage[italian]{babel}
\usepackage{comment}
\usepackage{tabularx}
\usepackage{longtable}
\usepackage{float}
\usepackage{amsmath}

\linespread{1.2}
\captionsetup[table]{labelformat=empty}
\geometry{headsep=1.5cm}

\renewcommand{\contentsname}{Indice}
\renewcommand\familydefault{\sfdefault}

\begin{document}

\newgeometry{left=2cm,right=2cm,bottom=2.1cm,top=2.1cm}
\begin{titlepage}
    \vspace*{.5cm}

    \vspace{2cm}
    {
        \centering
        {\bfseries\huge \Title\par}
        \bigbreak
        {\bfseries\large \Author\par}
        \bigbreak
        {\Version\par}
        \vfill

        \begin{tikzpicture}[remember picture,overlay]

            \fill[blue!20!black] (current page.south west) -- ++(10cm,0) -- ++(-10cm,15cm) -- cycle;
            \fill[orange] (current page.south east) -- ++(-18cm,0) -- ++(21.6cm,22cm) -- cycle;

            \clip (0,-3cm) circle (2.5cm) node (current page.center) {\includegraphics[width=5cm]{logo.jpg}};
        \end{tikzpicture}

    }

    \vfill

\end{titlepage}

\restoregeometry






















\newpage

%---------------header------------------%
\pagestyle{fancy}
\fancyhead{} % pulizia degli header
\lhead{%
\begin{tikzpicture}
    \clip (0,0) circle (0.5cm);
    \node at (0,0) {\includegraphics[width=1cm]{logo}};
\end{tikzpicture}%
}
\chead{\vspace{\fill}\Title\vspace{\fill}}
\rhead{\vspace{\fill}\Version\vspace{\fill}}
%quad è una spaziatura
%---------------------------------------%

\begin{table}[!h]
	\caption{Versioni}
	\footnotesize
	\begin{center}
		\begin{tabular}{ l l l l l }
			\hline \\[-2ex]
			Ver. & Data       & Redattore      & Verificatore      & Descrizione     \\
			\\[-2ex] \hline \\[-1.5ex]
			1.0  & 2024-04-24 & Leonardo Baldo & Antonio Benetazzo & Stesura verbale \\
			\\[-1.5ex] \hline
		\end{tabular}
	\end{center}
\end{table}

\newpage

\tableofcontents

\newpage

\section{Dettagli della riunione}


\textbf{Sede della riunione}: Piattaforma Discord\\
\textbf{Orario di inizio}: 15:00\\
\textbf{Orario di fine}: 16:00\\


\begin{flushleft}
	\begin{table}[!h]
	\begin{tabular}{ |l|l|l| } 
		\hline
		\textbf{Partecipante} & \textbf{Ruolo}       & \textbf{Presenza} \\
		\hline 
		Antonio Benetazzo     & Verificatore         & Presente          \\
		Davide Malgarise      &                      & Presente          \\
		Elena Ferro           &                      & Presente          \\
		Leonardo Baldo        & Redattore            & Presente          \\
		Matteo Tiozzo         &                      & Presente          \\
		Raul Seganfreddo      & Amministratore       & Presente          \\
		Valerio Occhinegro    &                      & Presente          \\
		\hline
	\end{tabular}
	\end{table}
\end{flushleft}

\section{Ordine del giorno}
\begin{itemize}
	\item Cosa è stato fatto fino ad ora
	\item Cosa rimane da fare
	\item Varie ed eventuali
	\item Decisioni prese e conclusioni
\end{itemize}

\newpage

\section{Verbale}

\subsection{Cosa è stato fatto fino ad ora}
Abbiamo fatto un riepilogo di quanto completato fino ad ora.
\begin{itemize}
	\item \textbf{Verbali} \\
	Redatto, revisionato e approvato il verbale della riunione esterna del 2024-04-22.
	\item \textbf{Avanzamento Documentazione} \\
	Abbiamo effettuato una prima stesura dell'\textit{Analisi dei Requisiti} mentre prosegue la stesura del \textit{Piano di Progetto}, del \textit{Piano di Qualifica}, delle \textit{Norme di Progetto}, e del 
	\textit{Glossario}.
	\item \textbf{Richieste dell'azienda} \\
	Abbiamo iniziato lo studio e l'analisi dei prossimi strumenti da utilizzare, ovvero \textit{ClickHouse} e \textit{Grafana} al fine di migliorare la persistenza dati implementando delle tabelle di aggregazione dei dati e al fine di integrare l'ultimo strumento dello stack tecnologico previsto. \\
	Prodotto inoltre il documento di analisi e confronto tra \textit{Apache Kafka} e \textit{Redpanda} per valutare in maniera più informata quale strumento sia opportuno utilizzare.
\end{itemize}

\subsection{Cosa rimane da fare}
Per questo sprint ci sono ancora alcune attività da completare:
\begin{itemize}
	\item finire la stesura delle \textit{Norme di Progetto};
	\item avanzamento e aggiornamento di \textit{Piano di Progetto}, \textit{Piano di Qualifica}, \textit{Analisi dei Requisiti} e \textit{Glossario};
	\item effettuare la stesura del verbale della riunione in corso.
\end{itemize}

\subsection{Varie ed eventuali}
Riguardo al documento di analisi e confronto tra \textit{Apache Kafka} e \textit{Redpanda}, l'azienda ha deciso che faremo delle considerazioni durante il prossimo SAL, previsto per il giorno 2024-05-06. In questo modo potremo parlarne assieme e valutare ulteriormente il nostro punto di vista prima di prendere una decisione definitiva.

\subsection{Decisioni prese e conclusioni}
Concludiamo la riunione impegnandoci a completare le attività rimanenti entro il termine del secondo sprint e ci diamo appuntamento per il prossimo incontro interno previsto eccezionalmente per il 2024-04-29, visto che la data originariamente prevista, 2024-05-01, è festa e molti componenti non sarebbero disponibili.

\end{document}